\chapter{General interface}
\label{chap:interface}

\section{Overview}

The main window of AirfoilTools is organised into four
areas (\figref{01_main_window.png}):

\begin{enumerate}
    \item the \textbf{menu bar} (\menu{File}, \menu{Edit},
        \menu{View}, \menu{Options}, \menu{Help});
    \item the \textbf{tab selector} (\menu{Airfoils},
        \menu{XFoil Settings}, \menu{Results},
        \menu{Cp / Boundary layer});
    \item the \textbf{work area}, which changes depending on the
        selected tab;
    \item the \textbf{status bar}, at the bottom of the window, where
        information messages and simulation progress are
        displayed.
\end{enumerate}

\section{Tooltips}

Each control (checkbox, input field, button, menu action)
has a \textbf{tooltip} that appears when the mouse
cursor hovers over it for a moment. These tooltips describe the role
of the control and the typical values.

\begin{astuce}
Hover over a control whose function is not obvious:
the tooltip appears after about a second.
For menu actions, the help is displayed at the bottom of
the window in the status bar when hovering over the
menu entry.
\end{astuce}

\section{\menu{File} menu}

\begin{tabularx}{\linewidth}{l l X}
    \toprule
    \textbf{Entry} & \textbf{Shortcut} & \textbf{Action} \\
    \midrule
    Open current airfoil\dots
        & \touche{Ctrl}+\touche{O}
        & Loads an airfoil into the current position (blue).
          Accepted formats: \chemin{.dat} (Selig/Lednicer),
          \chemin{.bspl} (AirfoilTools spline), \chemin{.bez}
          (AirfoilTools Bézier), \chemin{.csv}.\\
    Open reference airfoil\dots
        & \touche{Ctrl}+\touche{Shift}+\touche{O}
        & Loads an airfoil into the reference position (red).\\
    From the UIUC database / Current airfoil\dots
        & ---
        & Opens a search dialog in the online UIUC (Selig)
          database and downloads the selected airfoil.
          See section~\ref{sec:uiuc}, page~\pageref{sec:uiuc}.\\
    From the UIUC database / Reference airfoil\dots
        & ---
        & Same for the reference airfoil.\\
    NACA airfoil / Current airfoil\dots
        & ---
        & Generates a NACA airfoil (4 or 5 digits) as the current airfoil.
          Prompts for the digits and then for the number of points. See
          section~\ref{sec:naca}, page~\pageref{sec:naca}.\\
    NACA airfoil / Reference airfoil\dots
        & ---
        & Same for the reference airfoil.\\
    Save airfoil\dots
        & \touche{Ctrl}+\touche{S}
        & Saves the current airfoil. A first dialog lets you
          save the \emph{full airfoil}, the \emph{upper surface}
          only or the \emph{lower surface} only; the format is then
          selected from the filter of the file dialog.\\
    Save airfoil with flap\dots
        & ---
        & Saves the airfoil with the flap deflected (green), in the same
          way as the current airfoil. Requires the flap to be
          active. See section~\ref{sec:flap}, page~\pageref{sec:flap}.\\
    Open project\dots
        & \touche{Ctrl}+\touche{Shift}+\touche{P}
        & Opens an AirfoilTools project (\chemin{.aftproj}):
          restores the current and reference airfoils as well as
          the tracing image and its transformation. See
          chapter~\ref{chap:formats}.\\
    Save project\dots
        & \touche{Ctrl}+\touche{Alt}+\touche{S}
        & Saves the whole session (airfoils + tracing
          image) to a \chemin{.aftproj} file.\\
    Load tracing image\dots
        & ---
        & Loads an image (jpg, png, tif, bmp) in the background to
          trace an airfoil. See section~\ref{sec:calque},
          page~\pageref{sec:calque}.\\
    Remove tracing image
        & ---
        & Removes the tracing image from the current session.\\
    Quit
        & \touche{Ctrl}+\touche{Q}
        & Closes the application.\\
    \bottomrule
\end{tabularx}

\section{\menu{Edit} menu}

\begin{tabularx}{\linewidth}{l l X}
    \toprule
    \textbf{Entry} & \textbf{Shortcut} & \textbf{Action} \\
    \midrule
    Convert to Spline
        & \touche{Ctrl}+\touche{B}
        & Converts the current airfoil (discrete points mode)
          to multi-segment spline mode. See the section
          \emph{Conversion to spline}, page~\pageref{sec:convertspline}.\\
    Sampling / Current airfoil\dots
        & ---
        & Changes the number of sampling points of the
          current airfoil. Available only in Spline mode.\\
    Sampling / Reference airfoil\dots
        & ---
        & Same for the reference airfoil.\\
    \bottomrule
\end{tabularx}

\section{\menu{View} menu}

\begin{tabularx}{\linewidth}{l l X}
    \toprule
    \textbf{Entry} & \textbf{Shortcut} & \textbf{Action} \\
    \midrule
    Fit zoom
        & \touche{Ctrl}+\touche{0}
        & Reframes the \menu{Airfoils} canvas view on all
          visible airfoils.\\
    Layout / 1$\times$1, 1$\times$2, \dots
        & ---
        & Configures the grid of the \menu{Results} tab
          (rows $\times$ columns).\\
    \bottomrule
\end{tabularx}

\section{\menu{Options} menu}

\begin{tabularx}{\linewidth}{l l X}
    \toprule
    \textbf{Entry} & \textbf{Shortcut} & \textbf{Action} \\
    \midrule
    Deviation / Vertical (thickness)
        & ---
        & Computes the deviation between airfoils as a
          thickness difference measured \textbf{vertically} (along the $z$ axis),
          at constant abscissa. Default mode.\\
    Deviation / Normal (perpendicular)
        & ---
        & Computes the deviation \textbf{perpendicularly} to the
          surface of the current airfoil (distance along the normal).
          See section~\ref{sec:deviation},
          page~\pageref{sec:deviation}.\\
    \bottomrule
\end{tabularx}

\begin{noteinfo}
The two deviation modes are \textbf{exclusive}: choosing one
automatically unchecks the other. The change is reflected
immediately on the plot of the \menu{Airfoils} tab.
\end{noteinfo}

\section{\menu{Help} menu}

\begin{tabularx}{\linewidth}{l l X}
    \toprule
    \textbf{Entry} & \textbf{Shortcut} & \textbf{Action} \\
    \midrule
    User manual
        & \touche{F1}
        & Opens this manual (PDF file) with the system's default
          viewer. The PDF is embedded in the
          \chemin{.exe} distribution and remains accessible even
          offline.\\
    About\dots
        & ---
        & Displays version information and the author.\\
    \bottomrule
\end{tabularx}

\begin{noteinfo}
In development mode (launched via \chemin{run\_gui.py}), the
manual is read from \chemin{docs/manuel/manuel.pdf} in
the project tree. In PyInstaller executable mode, it is
extracted from the internal bundle (folder
\chemin{\_internal/docs/manuel.pdf} next to the exe).
\end{noteinfo}

\section{Tabs}

The four tabs are independent: switching tabs does not lose
any data. The status bar can display contextual messages
depending on the active tab. The \menu{Cp / Boundary layer} tab
(chapter~\ref{chap:cp}) provides a detailed, \emph{XFoil}-style view
of a single operating point.

\subsection*{Stale results indicator}

The \menu{Results} tab has a special mechanism:
when an airfoil is modified after the simulation has been run,
a \textbf{yellow banner} appears at the top of the tab to indicate
that the displayed results no longer match the
current airfoils. It is then recommended to re-run the simulation.

\begin{noteinfo}
The staleness banner does not block consultation of the
results: it is merely a visual indicator. The
banner is removed automatically when a new
simulation is launched.
\end{noteinfo}
